% paper5_amendments.tex -- Paper 5. Typeset alone with this file, or inside
% the volume through main.tex.
\documentclass[../main.tex]{subfiles}
\begin{document}
\begin{paper}{Organic Amendments and the Stabilisation of Soil Organic Carbon}{Organic amendments and soil carbon}
\label{paper:amendments}

\begin{paperabstract}
Adding organic materials to soil raises its carbon content, but only part of
the added carbon remains in the soil for long. This paper compares three
organic amendments, applied at equal rates of carbon, with an unamended
control at three sites, and follows topsoil organic carbon over the course of
the trial. The fraction of added carbon retained in the soil is calculated
for each amendment, and the retention is related to soil texture through the
clay content of each site. The data are synthetic and illustrate how carbon
retention is calculated and how it can be interpreted in terms of the
mechanisms that protect organic matter from decomposition.
\end{paperabstract}
\paperkeywords{organic amendments; manure; compost; soil organic carbon; carbon retention; clay}

\section{Introduction}

Organic amendments such as animal manure, compost and crop residues supply
nutrients and add carbon to the soil. They are a core component of
integrated soil fertility management, in which they are combined with
fertiliser to raise the efficiency with which nutrients are used
\parencite{vanlauwe2010}. Beyond their nutrient value they feed the soil
organisms on which many soil functions depend \parencite{doran2000}.

Most of the carbon in a fresh amendment is decomposed within a few seasons
and returned to the atmosphere as carbon dioxide. The fraction that persists
depends on how the remaining organic matter is protected. A review of the
mechanisms concluded that the chemical resistance of organic compounds to
decay is important only in the early stages of decomposition, and that
long-term persistence depends instead on the physical inaccessibility of
organic matter to decomposers and on its association with mineral surfaces
\parencite{vonlutzow2006}.

Because organic matter is protected by association with silt and clay
particles and by its occlusion within aggregates, soils differ in how much
carbon they can stabilise, and the protective capacity of a soil sets an
upper limit to the carbon it can hold in protected forms \parencite{six2002}.
Soils with more clay would therefore be expected to retain a larger share of
added carbon, while soils close to that upper limit would retain less of each
new addition. These ideas have practical weight, because the carbon that
agricultural soils could regain is considerable \parencite{lal2004}.

The aims of this paper are to compare the gain in topsoil organic carbon under
three amendments with that of an unamended control, to calculate the fraction
of the added carbon retained under each amendment, and to examine whether
retention is greater at sites with higher clay content.

\input{\subfix{methods/paper5_methods}}

\end{paper}
\end{document}
